\paragraph{A CPT constraint on the unchanged charged fields and clock.}
For a real, even, unit scalar Schwartz packet $\phi$, average the signed
addition/removal measure over all sixteen continuum field components.
Writing this probability measure as $\nu_\phi$, spin trace cancels the
Weyl first moment and its kinetic cross term. The exact finite-spacing
sum rules, $r_a\to1/2$, and weak lower semicontinuity give
\[
 \int E\,d\nu_\phi(E)=\mu=U_0/2,\qquad
 \int E^2\,d\nu_\phi(E)\le C_\phi
 =v^2\|\nabla\phi\|^2+U_0^2/2.
\]
Only the second-moment inequality passes to the limit. Negative response
frequency denotes removal; the two charge-tagged excitation energies
are nonnegative with respect to the original positive GNS generator $L$.

Suppose a proposed physical identification supplies an antiunitary
$\Theta$ preserving the same ground and $L$, and mapping the sixteen
smeared annihilators $A_s$ into their adjoints by a unitary matrix $Q$:
\[
 \Theta\Omega=\Omega,\qquad \Theta L\Theta^{-1}=L,\qquad
 \Theta A_s\Theta^{-1}=\sum_t Q_{st}A_t^*.
\]
The real even packet is preserved by the spatial inversion and conjugation
in the usual CPT action. For every spectral set $B$, antiunitarity and
unitarity of $Q$ give
\[
 \sum_s\|P_L(B)A_s^*\Omega\|^2
 =\sum_s\|P_L(B)A_s\Omega\|^2.
\]
Thus the averaged addition and removal measures agree, $\nu_\phi$ is
reflection symmetric, and its first moment is zero. This precise
same-field, same-ground, same-clock CPT identification is therefore
impossible for $U_0>0$. No CPT theorem is applied without its hypotheses:
this does not exclude Lorentz covariance alone or other field dictionaries.

With $\nu_\phi^-$ the reflected measure and total variation defined as
$\sup_B|\nu_\phi(B)-\nu_\phi^-(B)|$, Cauchy--Schwarz proves
\[
 \operatorname{TV}(\nu_\phi,\nu_\phi^-)
 \ge\frac{\mu^2}{C_\phi},\qquad
 W_1(\nu_\phi,\nu_\phi^-)\ge2\mu.
\]
There is a bounded odd witness: for $R=C_\phi/(2\mu)$, clipping $E$ to
$[-R,R]$ gives $f_R$ with $|E-f_R(E)|\le E^2/(4R)$ and
$\int f_R\,d\nu_\phi\ge\mu/2$. Gaussian packets of width $\ell$
have $C_\phi=3v^2/(2\ell^2)+U_0^2/2$. At $v=1/8$, $U_0=1$ and
$\ell=1$, the total-variation lower bound is $32/67$; broad packets make
it tend to $1/2$. These are rigorous bounds, not computed infinite-ground
spectra or executed native measurements.

A scalar charged-field rephasing changes the generator to
$L-\kappa Q_{\rm charge}$ and the signed first moment to $\mu-\kappa$.
Consequently $\kappa=U_0/2$ is the unique necessary centering for this
CPT condition. It is a change of charged dynamics, not an additive
energy-origin change or a rest mass. Positivity on the old ground is
not automatic: the empty two-mode ground of $H=n_0+3n_1$ has a centered
response of zero mean at $\kappa=2$, while $H-2N$ has eigenvalues
$0,-1,1,0$. The new reference and its ground property require proof.

The free $U(a)\to0$ control has physical addition and removal energies
$v|p|$, zero rest masses, and four Weyl chiralities of each sign. A
massive Dirac control, $H_D=Z\otimes(\sigma\cdot p)+mX\otimes I$,
has $H_D^2=(|p|^2+m^2)I$ and zero trace: its mass couples chiralities
and does not create the scalar signed-energy shift. Exact full-CAR
counterexamples and the conditional proof are retained in
\texttt{model/continuum\_cpt\_constraint}; no physical species, interacting
Lorentz completion or Einstein premise is thereby supplied.
